As generative artificial intelligence enters multimodal content production, official tourism accounts increasingly use AI-generated images, virtual scenes, and surreal visual expressions in short videos. This study examines tourism short videos published by provincial tourism accounts on Douyin. After screening and verification, 722 deeply reviewed samples are retained for formal analysis. Content analysis, non-parametric tests, multiple regression, robustness checks, and supplementary predictive analysis are used to evaluate the factors influencing user engagement. Public engagement indicators, including likes, comments, favorites, and shares, are used to measure user engagement.

The results show that engagement is highly skewed. Although several content variables show significant differences in univariate tests, multiple regression results indicate that AI technical quality, video style, video duration, and days since collection remain more stable predictors after controlling for other factors. Compared with low-quality AI output, moderate-quality and high-quality AI output significantly improve communication effectiveness. Surreal video style also has a significant positive effect compared with documentary or realistic style. Video duration has a negative effect, while days since collection has a positive effect on accumulated engagement. Machine learning results further show that AI technical quality, video duration, and video style have high predictive importance.

This study suggests that AI tourism short video marketing should not focus merely on whether AI is used or how deeply AI is involved. Instead, practitioners should pay more attention to the perceived quality, visual style, and viewing experience of AI-generated content.
